% Appendix --- supplementary material. Author: Stylianos Giannoulis (AUTH).

\section{Implementation Fidelity Audit}\label{app:fidelity}

Every detector implementation was audited against its published description and,
where a repository exists, against its official code, before any experiment was run.
Of the twenty-three candidate detectors, five matched their specification up to
cosmetic, metric-preserving differences (MSP, ODIN, Mahalanobis/MDS, DFM-PCA, SRS);
ten exhibited moderate or major deviations (DriftLens, TD-IVDM, DiMMAD, DEEDEE,
CODiT, M2N2, DIVERSIFY, InvAD, CatSight, AE-ADWIN-LSTM); and eight were \emph{not}
faithful reproductions (ReAct, DICE, GradNorm, SCALE, Outlier Exposure, DivOE,
DiffAD, DiverseMix). \Cref{tab:audit} summarises the verdicts by detector.

\begin{table}[h]
\centering
\caption{Faithfulness audit verdicts for all twenty-three candidate detectors. The
verdict states the severity of the discrepancy found between the implementation and
the source paper (and official code where available). ``Minor'' denotes
cosmetic/monotone differences that do not change threshold-free metrics; ``Moderate''
and ``Major'' denote discrepancies that can change rankings; ``Not verified'' denotes
a discrepancy that materially changes the OOD score.}
\label{tab:audit}
\small
\begin{tabular}{lll}
\toprule
Detector & Cluster & Verdict \\
\midrule
MSP & post-hoc softmax & Verified, minor deviations \\
ODIN & post-hoc softmax & Verified, minor deviations \\
Mahalanobis (MDS) & feature distance & Verified, minor deviations \\
DFM-PCA & feature density & Verified, minor deviations \\
SRS & ts-specific & Verified, minor deviation \\
DriftLens & drift/embedding & Verified, moderate deviation \\
TD-IVDM & concept drift & Verified, moderate deviation \\
DiMMAD & feature distance & Verified, moderate deviation \\
DEEDEE & ts-specific & Verified, major deviation \\
CODiT & ts-specific & Verified, moderate deviation (orientation) \\
M2N2 & test-time adaptation & Verified, moderate deviation (order-dependent) \\
DIVERSIFY & ts-specific & Verified, major deviation \\
InvAD & ts-specific & Verified, major deviation \\
CatSight & concept drift & Verified, moderate deviation \\
AE-ADWIN-LSTM & concept drift & Verified, major deviation \\
ReAct & activation shaping & Not verified (score family) \\
DICE & activation shaping & Not verified (abs-sum + MSP) \\
GradNorm & gradient & Not verified (wrong statistic) \\
SCALE & activation shaping & Not verified (wrong layer) \\
Outlier Exposure & augmented training & Not verified (no outlier training) \\
DivOE & augmented training & Not verified (no synthesis/training) \\
DiffAD & reconstruction & Not verified (input-independent reverse) \\
DiverseMix & augmented training & Not verified (score inverted) \\
\bottomrule
\end{tabular}
\end{table}

\paragraph{The six corrected variants.}
Six discrepancies admitted a paper-faithful correction, implemented and validated as
\texttt{\_enh} variants that enter the main sweep as an original-versus-enhanced
ablation (\cref{tab:ablation_deltas}). \texttt{gradnorm\_enh} computes the $L_1$ norm
of the last-layer-weight gradient of the KL-to-uniform objective and negates it
(matching the published confidence orientation). \texttt{dice\_enh} sums the signed
top-$k$ contributions and computes energy. \texttt{react\_enh} computes energy after
activation clipping. \texttt{scale\_enh} applies activation scaling at the
penultimate layer. \texttt{dimmad\_enh} drops the binary set metrics (Hamming,
Jaccard, Dice) that are ill-defined on continuous features.
\texttt{diversemix\_enh} negates the score to match the training objective; on the
synthetic validation task this raised AUROC from \num{0.188} to \num{0.812}, yet on
real streaming data the same correction lowered mean AUROC from \num{0.633} to
\num{0.367} --- the regime-dependent orientation result discussed in the main text.

\paragraph{The feature-space relabelling.}
The audit also established that the per-sample implementations of DriftLens
(PCA-Mahalanobis), TD-IVDM (KDE density), DiMMAD (distance ensemble), DFM-PCA (PCA
reconstruction), and Mahalanobis are all variants of distance/density estimation in
the frozen feature space. The robust winning signal is therefore feature-space
distance/density surfacing under several names, which reframes the headline result.

\paragraph{The systemic finding.}
At least seven detectors --- GradNorm, DiverseMix, DiffAD, AE-ADWIN-LSTM, CatSight,
CODiT, and the M2N2 EMA path --- contained a score negation or orientation choice
rationalised post-hoc by an observation on small training data rather than following
the principled score from the paper. Several score below chance on the real sweep
precisely because the fitted orientation does not generalise. Together with the
absent OE/DivOE training, GradNorm's wrong statistic, DiffAD's input-independent
reverse process, and the relabelling above, this means an unaudited ranking would
substantially reflect implementation artifacts. This is the central methodological
finding and the rationale for the verify-before-experiment gate.

\section{Original-versus-Corrected Ablation}\label{app:ablation}

\Cref{tab:ablation_deltas} (reproduced in the main text) reports the mean-AUROC
effect of each correction. A Friedman test over the corrected-variant comparison
remains significant both before correction ($\chi^2 = 8.687$, $p = 0.0338$) and after
($\chi^2 = 16.722$, $p = 0.0008$), and the top-five ranking --- SRS, Mahalanobis,
DriftLens, TD-IVDM, DFM-PCA --- is identical under both. The conclusion is two-part:
the implementations contained real, conclusion-relevant bugs, and the headline
finding is robust to fixing them.

\begin{figure}[h]
\centering
\includegraphics[width=0.85\textwidth]{ablation_delta.png}
\caption{Change in mean AUROC from the original implementation to the paper-faithful
corrected (\texttt{\_enh}) variant, for the six corrected detectors. Positive bars
denote improvement. The gradient-norm repair yields the largest gain ($+0.122$) yet
still lands below chance, while the diversification correction reverses on real data
($-0.266$); no correction lifts a post-hoc detector into competitiveness with the
feature-space methods.}
\label{fig:app_ablation}
\end{figure}

\section{ROC Curves and Score-Distribution Case Study}\label{app:roc}

\begin{figure}[h]
\centering
\includegraphics[width=0.95\textwidth]{appendix_roc.png}
\caption{Receiver-operating-characteristic curves for three leading feature/drift
detectors and three post-hoc softmax detectors on a representative
global-normalised dataset. The softmax curves fall below the diagonal, the signature
of score inversion, while the feature-space curves hug the top-left corner.}
\label{fig:app_roc}
\end{figure}

\begin{figure}[h]
\centering
\includegraphics[width=0.9\textwidth]{appendix_violin_case.png}
\caption{Distribution of out-of-distribution scores for in-distribution versus
out-of-distribution test windows on a global-normalised dataset. The feature-space
detector (left) assigns systematically higher scores to OOD windows; the softmax
detector (right) assigns them \emph{lower} scores, inverting the decision and
producing an AUROC below chance --- the mechanism analysed in the main text.}
\label{fig:app_violin}
\end{figure}

\section{Reproducibility}\label{app:repro}

All experiments fix the random seed to \num{42} and run on CPU under Python~3.13 with
PyTorch, NumPy, SciPy, and scikit-learn. The backbone is a 1-D ResNet trained for
\num{40} epochs with cross-entropy over temporal pseudo-classes. The benchmark
comprises \num{36} datasets (\num{18} TSB-StreamingAD-U and \num{18}
TSB-StreamingAD-M, six per DRIFT/OOD/STABLE cell), \num{22} detectors, and \num{772}
completed runs. Per-run configurations and metrics are stored alongside the aggregate
results matrix, and the tables and figures in this manuscript are regenerated
directly from it. The reduced epoch budget and the stratified subsample of \num{36}
of the \num{600} available files are recorded explicitly as scoping decisions; SRS is
excluded on TSB-StreamingAD-M because it hangs on multivariate input, and Outlier
Exposure and DivOE are reported as energy baselines because faithful
auxiliary-outlier training is not available.
